% Structure du corpus. Le déséquilibre est celui de la loi elle-même : le
% Livre II concentre à lui seul près de la moitié des articles, ce qui
% explique que le jeu de référence y puise autant.
% Répartition des articles du corpus par Livre
% Généré par rapport/build_data.py — NE PAS ÉDITER.
\newcommand{\CorpusLabels}{{Livre II},{Livre I},{Livre III},{Livre IV},{Livre VI},{Livre V},{Livre Prélim.},{Livre VII}}
\newcommand{\CorpusCoords}{(261,0) (122,1) (79,2) (55,3) (37,4) (19,5) (12,6) (4,7)}
\newcommand{\CorpusMax}{261}

\begin{tikzpicture}

\begin{axis}[
  name=livres,
  barresH,
  height=5.4cm, width=0.55\linewidth,
  xmin=0, xmax=\CorpusMax*1.28,
  ytick={0,...,7},
  % Libellés en clair : pgfplots ne développe pas une macro donnée à
  % « yticklabels ». L'ordre suit celui des barres, par effectif décroissant.
  yticklabels={{Livre II},{Livre I},{Livre III},{Livre IV},
               {Livre VI},{Livre V},{Livre prélim.},{Livre VII}},
  yticklabel style={font=\sffamily\scriptsize, text=ink, align=right},
  xlabel={nombre d'articles},
  nodes near coords,
  bar width=9pt,
  title={Répartition par Livre},
  title style={at={(0,1.06)}, anchor=south west},
]
\addplot[draw=inptblue, fill=inptblue, fill opacity=0.85]
  coordinates {\CorpusCoords};
\end{axis}

% Décomposition du Livre II — le seul dont la taille justifie qu'on l'ouvre.
\begin{axis}[
  barresH,
  at={($(livres.north east)+(1.7cm,0)$)}, anchor=north west,
  height=5.4cm, width=0.32\linewidth,
  xmin=0, xmax=118,
  ytick={0,...,4},
  yticklabels={{Durée du travail},{Hygiène et sécurité},{Salaire},
               {Mineur et femme},{Dispositions générales}},
  yticklabel style={font=\sffamily\scriptsize, text=ink, align=right},
  xlabel={articles},
  nodes near coords,
  bar width=8pt,
  title={Livre II, par Titre},
  title style={at={(0,1.06)}, anchor=south west},
]
\addplot[draw=inptblue, fill=inptblue, fill opacity=0.45]
  coordinates {(97,0) (64,1) (51,2) (41,3) (8,4)};
\end{axis}

% Le commentaire sur les métadonnées est passé en légende : posé dans la
% figure, il débordait de la justification.

\end{tikzpicture}
